\chapter{Related Work}
\label{chap:ch4}

In the following chapter, we review the evolution of methodologies in the literature designed to solve the 3D reassembly problem.

\section{Historical Evolution of 3D Shape Reassembly Methodologies}

Before the advent of deep learning, the field relied on \textbf{Traditional Geometry-Based Fracture Assembly}. These early methods utilized hand-crafted mathematical invariants and surface characteristics to align broken fragments. For example, Huang et al. \cite{huang2006reassembling} proposed reassembling fractured objects via geometric matching using integral invariants to analyze fracture surfaces. In archaeological contexts, Toler et al. \cite{toler2010multi}. utilized multi-feature matching, combining color, normal maps, and geometric features to reassemble fresco fragments, while Funkhouser et al. \cite{funkhouser2011learning} applied machine learning classifiers to evaluate these hand-crafted features and learn how to match the fragments effectively (Figure~\ref{fig:feature_descriptors}).

\begin{figure}[h]
\captionsetup[subfigure]{labelformat=empty}
    \centering

    \begin{subfigure}[t]{0.15\textwidth}
        \centering
        \includegraphics[height=2.3cm]{figures/brush_strokes.png}
        \caption{\tiny{Brush Strokes}}
    \end{subfigure}
    \hfill
    \begin{subfigure}[t]{0.15\textwidth}
        \centering
        \includegraphics[height=2.3cm]{figures/string_impressions.png}
        \caption{\tiny{String Impressions}}
    \end{subfigure}
    \hfill
    \begin{subfigure}[t]{0.15\textwidth}
        \centering
        \includegraphics[height=2.3cm]{figures/curvature.png}
        \caption{\tiny{Curvature}}
    \end{subfigure}
    \hfill
    \begin{subfigure}[t]{0.15\textwidth}
        \centering
        \includegraphics[height=2.3cm]{figures/erosion.png}
        \caption{\tiny{Erosion}}
    \end{subfigure}
    \hfill
    \begin{subfigure}[t]{0.15\textwidth}
        \centering
        \includegraphics[height=2.3cm]{figures/decorative_imagery.png}
        \caption{\tiny{Decorative Imagery}}
    \end{subfigure}
    \hfill
    \begin{subfigure}[t]{0.15\textwidth}
        \centering
        \includegraphics[height=2.3cm]{figures/texture.png}
        \caption{\tiny{Texture}}
    \end{subfigure}

    \caption{Examples of hand-crafted surface feature descriptors used in traditional fragment matching methods \cite{toler2010multi}. Each descriptor captures a different physical property of the fragment surface, from brush stroke patterns to curvature profiles.}
    \label{fig:feature_descriptors}
\end{figure}

\begin{wrapfigure}{r}{0.45\textwidth}
    \centering
    \vspace{-12pt}
    \includegraphics[width=0.43\textwidth]{figures/skull.png}
    \caption{Template-guided reassembly of a fractured skull \cite{yin2011automatic}.}
    \label{fig:skull_template}
    \vspace{-10pt}
\end{wrapfigure}
Because pure geometric matching struggles when fragments are severely damaged or missing, a sub-branch of research focused on \textbf{Template-Guided Assembly}. These methods rely on a complete reference model to guide the placement of fragments. A prominent example is the automatic assembly and completion framework proposed by Yin et al. \cite{yin2011automatic}, which roughly aligns fragmented skulls using a healthy template skull before refining the break-curve matches (Figure~\ref{fig:skull_template}).

With the rise of neural networks, the focus shifted towards \textbf{Deep Learning for Semantic Part Assembly}. Rather than physical fractures, these models focused on constructing complete objects from predefined, semantically meaningful parts (e.g., table legs or chair backs). Pioneering methods in this domain include Generative 3D Part Assembly via Dynamic Graph Learning (DGL), which utilizes iterative graph message passing to reason about spatial relationships and refine part poses. \cite{Huang2020DGL}. This was subsequently improved by RGL-Net, a recurrent graph learning framework \cite{harish2022rgl}. Other approaches used different constraints, such as predicting 3D part assembly guided by a single 2D image \cite{li2020learning}, or utilizing PQ-NET, a generative seq2seq network that sequentially predicts part poses to form 3D shapes \cite{wu2020pq}.

Although highly successful, these semantic methods rely heavily on prior structural knowledge and fail when applied to real-world objects shattered into meaningless fragments. This necessitated a transition toward \textbf{Deep Geometric Fracture Assembly}, which relies purely on the geometric cues of fracture surfaces. A notable method in the space is Neural Shape Mating (NSM) \cite{chen2022neural}, which leverages a Transformer network to predict poses and an adversarial shape prior to evaluate the visual plausibility of the mated geometry (Figure~\ref{fig:nsm}). However, NSM is strictly specialized for pairwise alignment (assembling exactly two pieces) and struggles to scale to complex multi-piece objects \cite{chen2022neural}.

\begin{figure}[h]
    \centering
    \includegraphics[width=1\linewidth]{figures/neural_shape_mating.png}
    \caption{The Neural Shape Mating pipeline \cite{chen2022neural}. A Transformer-based network predicts relative poses for two fragments, using an adversarial shape prior to evaluate the geometric plausibility of the assembled result.}
    \label{fig:nsm}
\end{figure}

A major turning point in the literature occurred with the release of the Breaking Bad dataset, a large-scale synthetic dataset that simulates how real-world objects break using physically-based fracture algorithms \cite{sellan2022breaking}. By removing semantic part consistency, this dataset shifted the direction of exploration definitively toward multi-piece fracture assembly. In this context, Jigsaw emerged as a historical milestone. It is widely recognized as the first learning-based framework specifically designed for the assembly of multiple, physically broken 3D pieces. By using a joint-learning pipeline with attention layers to simultaneously optimize point-wise surface segmentation and global multi-part matching, Jigsaw successfully bridged the gap between semantic-free geometric mating and multi-piece global alignment \cite{lu2023jigsaw}.

Consequently, Jigsaw established the standard baseline for subsequent research in this domain. Its pipeline consists of four stages: PointNet++ feature extraction with self-attention and cross-attention, binary fracture surface segmentation, global matching via primal-dual descriptors and Sinkhorn optimal transport, and RANSAC-based pose recovery with Shonan synchronization \cite{lu2023jigsaw}. The full pipeline is described in detail in Chapter~\ref{chap:ch5}. Although Jigsaw established a robust baseline for multi-part assembly, its architecture presented several critical bottlenecks. 

\section{New generation of SOTA}

In the following section we examine the Jigsaw framework, which, despite establishing a strong baseline for multi-fragment reassembly, introduces several structural and computational inefficiencies within its matching and alignment pipeline. These shortcomings reveal important bottlenecks related to scalability, geometric reasoning, and pose estimation accuracy. Based on this analysis, we then review how newer methods address these limitations through more efficient matching strategies and improved alignment mechanisms.

\subsection{The Efficiency Bottleneck}

First, relying on dense point-wise matching—where every single point on a fragment is compared to every point on another—is computationally expensive, making it difficult to scale to complex objects with many parts.

\paragraph{PHFormer: Multi-Fragment Assembly Using Proxy-Level Hybrid Transformer} \mbox{} \\
To address the computational bottleneck of dense point matching, the PHFormer \cite{cui2024phformer} model adopts a patch-based representation. Instead of operating on individual points, the input point cloud is partitioned into a set of local regions, referred to as proxies (Figure~\ref{fig:phformer}). These proxies are constructed by first selecting representative points using Farthest Point Sampling \cite{qi2017pointnetplusplus}, and then grouping their neighboring points to form local patches. For each proxy, an edge convolution operation is applied to extract a feature embedding that captures its local geometric structure. This transformation reduces the complexity of the matching problem while preserving essential spatial information.

\begin{figure}[H]
    \centering
    \includegraphics[width=1\linewidth]{figures/phformer.png}
    \caption{The PHFormer architecture \cite{cui2024phformer}. Point clouds are partitioned into proxy-level patches, which are processed through iterative intra- and inter-fragment attention layers before decomposing pose estimation into adjacency prediction and transformation regression.}
    \label{fig:phformer}
\end{figure}

PHFormer processes these proxy features using a two-level attention mechanism applied iteratively: \textbf{intra-fragment attention} builds a coherent representation of each fragment's internal geometry, while \textbf{inter-fragment attention} enables reasoning about spatial relationships across fragments.

Based on these enriched representations, PHFormer decomposes pose estimation into adjacency prediction, relative pose estimation, and absolute pose regression. Fragments form a graph whose edges are weighted by predicted adjacency probabilities, and message passing over this graph produces globally consistent rigid poses.

\paragraph{3D Geometric Shape Assembly via Efficient Point Cloud Matching} \mbox{} \\
A complementary approach to the efficiency problem is the Proxy Match Transformer (PMTR) \cite{lee20243d}, which replaces dense point-to-point comparisons with a sub-quadratic matching paradigm based on a shared proxy space (Figure~\ref{fig:pmtr}).

The pipeline begins with \textbf{multi-scale feature extraction}. A U-Net-style \cite{ronneberger2015u} backbone processes the raw point clouds of two fragments $X$ and $Y$ and produces feature maps at three spatial resolutions: coarse, medium, and fine. This allows the network to capture both the macro-structure of the fragments and the fine-grained details of the fracture edges.

The next stage is \textbf{coarse matching}. In standard geometric matching, evaluating the compatibility of two fragments requires computing a dense pairwise correlation matrix $F_X F_Y^\top$, which has $O(N \times M)$ quadratic memory complexity. Instead of this direct computation, PMTR introduces the Proxy Match Transform (PMT) layer. The key idea is to introduce a small, learnable Proxy Tensor $P$---effectively a virtual set of points whose dimensionality is much smaller than $F_X$ or $F_Y$. Each fragment independently projects its features onto this shared proxy. For a single attention head, the transform applies a second-order convolution (which, by the authors' lemma, is equivalent to a head of self-attention):
\begin{equation}
\label{eq:pmt}
PMT(F_X) = \sum_{h=1}^{N_h} A_X^{(h)} F_X P^{(h)\top} w_X^{(h)}
\end{equation}
where $A_X$ is a local attention matrix containing binary values that restricts computation to spatial neighborhoods, $w_X$ is a learnable scalar weight, and $P^\top$ contains the proxy features. This operation refines the initial features of each fragment while capturing the essence of their interaction, because the dot-product of the two independent outputs approximates the full pairwise convolution: $PMT(F_X) \cdot PMT(F_Y)^\top \approx \text{Conv}(F_X, F_Y)$---but without ever computing the quadratic correlation matrix. From these refined coarse features, PMTR computes a score matrix that identifies candidate mating regions, laying the foundation for more granular alignment.

\begin{figure}[H]
    \centering
    \includegraphics[width=1\linewidth]{figures/pmtr.png}
    \caption{The PMTR architecture \cite{lee20243d}. A coarse-to-fine matching hierarchy uses learned Proxy Tensors to establish sub-quadratic correspondences across multiple spatial resolutions, culminating in SVD-based rigid pose estimation.}
    \label{fig:pmtr}
\end{figure}

The pipeline then proceeds to \textbf{fine matching}. The same PMT layers are applied at medium and fine resolutions, with skip connections from the U-Net backbone progressively refining spatial precision. Fine-level features are grouped around their spatially proximate coarse matches, ensuring that the hierarchical refinement is guided by the coarse structure.

Finally, the resulting score matrix is transformed by a Sinkhorn Optimal Transport layer \cite{cuturi2013sinkhorn}, which enforces a doubly-stochastic constraint to produce a confidence matrix of one-to-one point correspondences. From these refined correspondences, a Singular Value Decomposition \cite{golub2013matrix} step computes the rigid 6-DoF transformation (rotation and translation) that assembles the fragments.

\subsection{The Attention Bottleneck}

Standard attention mechanisms are limited in assembly tasks because they rely mainly on feature similarity, without explicitly considering spatial constraints such as distances, angles, or relative orientations between fragments. These geometric relationships are essential for determining how parts fit together, yet vanilla cross-attention ignores them entirely. To address this limitation, recent methods aim to incorporate geometric information directly into the attention process.

\paragraph{Geometric Point Attention Transformer for 3D Shape Reassembly} \mbox{} \\
The Geometric Point Attention Transformer (GPAT) \cite{Li2023GPAT} addresses the attention bottleneck by explicitly encoding geometric relationships into the attention computation.

The architecture begins with a PointNet backbone that extracts local geometric features ($h_i^{local}$) from each fragment. These are pooled across all parts to form a global context feature ($h^{global}$), providing a structural shape prior. The local and global features are concatenated with recycled geometric embeddings---initialized to zero on the first pass---and processed through an MLP to produce node features ($h_i$). Simultaneously, pairwise cross-part features ($z_{ij}$) are constructed by concatenating the local features of fragments $i$ and $j$ together with recycled geometric pair features.

To iteratively refine these features, GPAT introduces a specialized attention mechanism composed of three branches:

\begin{itemize}
    \item \textbf{Part Attention} acts as a global context extractor. Standard dot-product attention over the node features ($h_i$) evaluates high-level interactions between all fragments, forming a fully connected graph where parts attend to one another.
    \item \textbf{Pair Attention} introduces explicit spatial constraints. It transforms the cross-part features ($z_{ij}$) into edge attention terms, encoding invariant geometric signals: the distances between the centers of mass of each pair of pieces and the triplet-wise dihedral angles between them, mapped into a continuous vector space via Radial Basis Functions (RBFs). This spatial bias is crucial for maintaining geometrically valid relative placements.
    \item \textbf{Point Attention} models 6-DoF pose interactions directly. Each part's node features are mapped to $N$ virtual 3D keypoints in Euclidean space. The currently estimated rigid poses ($T_i$, $T_j$) are applied to these query and key points, and alignment quality is measured via $\|T_i \circ q_i^m - T_j \circ k_j^m\|_2$.
\end{itemize}

The outputs of the three branches are fused into a single attention score: $a_{ij} = \text{softmax}(n_{ij} + e_{ij} - p_{ij})$. The Point Attention term ($p_{ij}$) is subtracted because it represents a spatial distance error: larger misalignment penalizes the attention score.

Rather than predicting absolute poses, the network outputs a relative update ($\delta R_i, \delta t_i$) at each iteration. GPAT then employs \textbf{Geometric Recycling}: the updated poses are physically applied to the point clouds, which are passed back through the PointNet backbone to re-extract spatially aware features for the next iteration. This progressive refinement is central to the architecture's effectiveness---the geometric descriptors become increasingly informative as the fragments approach their correct positions (Figure~\ref{fig:gpat}).

\begin{figure}[H]
    \centering
    \includegraphics[width=1\linewidth]{figures/gpat.png}
    \caption{Overview of the GPAT architecture \cite{Li2023GPAT}. The attention mechanism integrates feature similarity (Part Attention), geometric pair descriptors (Pair Attention), and point-level pose alignment (Point Attention). Geometric Recycling feeds updated poses back through the feature extractor for iterative refinement.}
    \label{fig:gpat}
\end{figure}

\subsection{The Blindness Bottleneck}

Classical reassembly methods rely primarily on local geometric cues, such as matching fracture surfaces. While effective in controlled settings, these approaches often fail when fragments are incomplete, eroded, or misaligned, as they lack an explicit representation of the global object structure.

\paragraph{Jigsaw++: Learning Global Shape Priors for Robust Reassembly} \mbox{} \\

Jigsaw++ \cite{lu2024jigsawpp} addresses this limitation by introducing a generative global shape prior. Rather than directly estimating fragment poses, the method first hallucinates a plausible complete object from a partial assembly, which then serves as a structural guide for the final geometric alignment.

The architecture operates in two phases (Figure~\ref{fig:jigsawpp}). In the first phase, the network is trained exclusively on intact, unbroken objects to learn what complete shapes look like. The 3D geometry is projected into multi-view 2D images and passed through LEAP \cite{jiang2024leap}, a pre-trained image-to-3D encoder built on DINOv2 vision features. This produces latent representations of perfect objects, forming the target distribution $\pi_1$. A \textbf{Rectified Flow} model then learns an Ordinary Differential Equation that transports samples from a standard Gaussian noise distribution $\pi_0$ along straight-line trajectories in latent space until they match $\pi_1$. At the end of this phase, the generative model can create complete 3D shapes from noise.

In the second phase, the actual broken object is introduced. The partial assembly is rendered and encoded through LEAP, producing a biased, incomplete latent representation $\hat{x}_0$. Because this represents a broken shape, it sits in a low-likelihood region of $\pi_0$. To initiate repair, Jigsaw++ applies a controlled perturbation via Langevin dynamics, pushing the broken latent into a better starting position $x_0$. The Rectified Flow model then performs a \textbf{retargeting} step: it fine-tunes the generative flow to transport this perturbed broken latent directly onto the corresponding complete latent $x_1$.

\begin{figure}[H]
    \centering
    \includegraphics[width=1\linewidth]{figures/jigsaw++.png}
    \caption{Overview of the Jigsaw++ pipeline \cite{lu2024jigsawpp}. Partial assemblies are rendered into multi-view images and encoded into a latent space. Rectified Flow transports the incomplete latent toward the manifold of valid complete shapes, providing a global structural prior for subsequent geometric alignment.}
    \label{fig:jigsawpp}
\end{figure}

The output of this process is not a final alignment but a coherent global shape reconstruction. This reconstructed shape provides strong geometric guidance for subsequent matching methods, effectively enabling the pipeline to reason about what the assembled object \emph{should} look like even when fragments are incomplete or severely misaligned.

\subsection{The Synthetic-to-Real Gap}

The majority of the previous models were trained and tested on the synthetic Breaking Bad dataset. Synthetic fractures lack the complex deterioration, erosion, and extraneous missing parts found in the real world archaeological or medical fractures.

\paragraph{GARF: Flow-Based 3D Reassembly on SE(3)} \mbox{} \\
Unlike prior methods such as PHFormer and PMTR, which formulate reassembly as a discrete matching or regression problem, GARF \cite{Li2025GARF} recasts the task as a \emph{continuous dynamical system} defined on the Lie group \(SE(3)\). Instead of directly predicting final poses, the model learns how fragments move over time from an initial noisy configuration to a coherent assembly.

The pipeline begins by extracting local geometric features from each fragment using a Point Transformer V3 (PTv3) backbone---a substantially more powerful feature extractor than the PointNet++ used in Jigsaw. A segmentation head identifies fracture surfaces, and Poisson disk sampling is applied to uniformly resample input point clouds, preventing the network from memorizing training-specific point density distributions. From these features, the network constructs two complementary embeddings: a \emph{shape embedding}, formed by concatenating the PTv3 features with positional encodings of point coordinates, surface normals, and fragment scale; and a \emph{pose embedding} that encodes the fragment's current spatial configuration. These two embeddings are summed and positionally encoded to produce the final representation.

This representation is processed through a deep Transformer. Intra-fragment self-attention, implemented via FlashAttention, refines each fragment's internal geometric understanding. Global cross-attention enables all fragments to exchange spatial context. A Feed-Forward Network processes the resulting features and predicts the rotational and translational velocities for the next integration step.

Each fragment pose \(T^i = (r^i, a^i) \in SE(3)\), consisting of a rotation \(r^i \in SO(3)\) and a translation \(a^i \in \mathbb{R}^3\), evolves according to a learned velocity field:
\begin{equation}
\label{eq:garf_flow}
\frac{d}{dt}\psi_t(T) = v_t(\psi_t(T)), \quad v_t \in \mathfrak{se}(3),
\end{equation}
where \(\mathfrak{se}(3)\) is the tangent space encoding instantaneous rigid motion. The initial configuration at \(t=0\) consists of uniformly random rotations and Gaussian-distributed translations; the target at \(t=1\) is the ground-truth assembly. The rotation and translation components are decoupled and updated independently. Given predicted angular velocity \(\hat{\omega}^i_t\) and translational velocity \(\hat{v}^i_t\), the pose evolves from time \(t\) to \(t + \Delta t\) as:
\begin{align}
\label{eq:garf_rot_update}
r^i_{t+\Delta t} &= \exp\bigl(\Delta t \cdot \hat{\omega}^i_t\bigr) \cdot r^i_t, \\
\label{eq:garf_trans_update}
a^i_{t+\Delta t} &= a^i_t + \Delta t \cdot \hat{v}^i_t,
\end{align}
where \(\exp(\cdot)\) denotes the exponential map from \(\mathfrak{so}(3)\) to \(SO(3)\). The network is trained via a flow-matching objective that aligns the predicted velocities with the ground-truth geodesic velocities connecting the noisy and target poses on \(SE(3)\).

A key challenge in this formulation is global ambiguity: the entire assembly can drift rigidly in space, making the learning signal ill-posed. GARF resolves this by randomly selecting a subset of \(k\) fragments as \emph{anchors} at each training iteration. Anchor fragments are frozen with zero velocity, and the network must assemble the remaining pieces relative to them. Because the anchor set changes every iteration, the model learns consistent relative geometric relationships rather than memorizing absolute positions.

At inference time, GARF employs a two-stage strategy (Figure~\ref{fig:garf}). A one-step pre-assembly forces the model to predict a coarse alignment in a single large step from \(t=0\) to \(t=1\), collapsing the search space. A multi-step ODE solver then refines this estimate through iterative integration. To bridge the gap to real-world fragments, GARF uses few-shot domain adaptation via LoRA (Low-Rank Adaptation): trainable adapters are injected only into the attention layers, allowing the model to adapt with as few as 5--10 real-world objects. The authors introduce the FRACTURA dataset---containing real-world fractures from ceramics, bones, eggshells, and lithics---as the first systematic benchmark for real-world fracture assembly.

\begin{figure}[H]
    \centering
    \includegraphics[width=1\linewidth]{figures/garf.png}
    \caption{Overview of the GARF pipeline \cite{Li2025GARF}. Fragment poses evolve along learned velocity fields on $SE(3)$, from an initial noisy configuration (left) to a coherent assembly (right) through a two-stage coarse-to-fine integration scheme.}
    \label{fig:garf}
\end{figure}

\subsection{Summary and Research Gap}

Recent deep learning methods have significantly improved the problem of 3D fractured object reassembly, especially with the introduction of frameworks such as Jigsaw, which combine attention mechanisms with global matching strategies. However, despite these advances, several limitations remain, particularly when dealing with ambiguous or incomplete fragments.

First, standard cross-attention relies mainly on learned feature similarity and does not explicitly consider geometric relationships between fragments, which can reduce robustness in complex multi-piece scenarios. Second, most architectures use relatively shallow attention, limiting their ability to iteratively refine feature representations. Finally, pairwise transformations are typically recovered using sparse correspondences through RANSAC, which does not fully utilize the dense geometric information available on fracture surfaces, potentially limiting alignment precision.

This thesis addresses these limitations by exploring targeted improvements within the Jigsaw pipeline. Instead of proposing a completely new architecture, we focus on understanding how specific modifications can influence the reconstruction process. In particular, we introduce a geometric bias into the attention mechanism, extend the model with a gated second attention pass for controlled refinement, and incorporate an ICP-based step to refine pairwise transformations before global alignment.

In the next chapter, we present the extended methodology and analyze how these modifications affect the overall reconstruction, with a focus on challenging scenarios such as incomplete or ambiguous fragment configurations.